\section{Physical caps and the residue pairing}
\label{sec:caps}
The raw real-cycle period is a useful holomorphic diagnostic, but it is
not yet a functional on physical vacua. This section constructs such
functionals and fixes their normalization. All norms use ordinary
Euclidean measure on $\C^2$, with $\norm{dz_j}^2=2$. We take Hermitian
inner products to be antilinear in the first slot.

\subsection{Compact representatives and long cylinders}
For a strongly elliptic polynomial $f$, set
\[
D_f=\bar\partial+df\wedge,\qquad
D_f^\dagger=2\sum_j
 \left(-\iota_{\bar j}\partial_{z_j}+\overline{f_j}\iota_j\right),
\qquad H_f=2\{D_f,D_f^\dagger\}.
\]
The realization is the closed Hilbert complex on square-integrable
forms. Its finite-dimensional harmonic space is denoted $\mathcal V_f$.
The polynomial models here satisfy the noncompact Hodge hypotheses;
the degenerate homogeneous endpoint is checked in
Section~\ref{sec:endpoints}. The relevant comparison theorem does not
require Morse critical points \cite{LiWen}.

Choose a compactly supported smooth cutoff $\chi$ equal to one near
every critical point. Off this set, define
\[
V_f=\frac{(df\wedge)^\dagger}{|df|^2},\qquad
N_f=[\bar\partial,V_f],\qquad
\mathcal R_f=V_f(1+N_f)^{-1}.
\]
Brackets here are graded. The inverse is a finite series by bidegree,
and $[D_f,\mathcal R_f]=1$. Consequently
\begin{equation}\label{caps:cutoff}
T_\chi=1-[D_f,(1-\chi)\mathcal R_f]
       =\chi+\bar\partial\chi\wedge\mathcal R_f,
\qquad [D_f,T_\chi]=0.
\end{equation}
For $\Omega=dz_1\wedge dz_2$ and a polynomial $p$, the form
$u_p=T_\chi(p\Omega)$ is compactly supported and closed. Its physical
cap is
\begin{equation}\label{caps:heat}
h_f(p)=P_fu_p=\lim_{t\to\infty}e^{-tH_f}u_p.
\end{equation}
This is a state preparation by a long Euclidean cylinder. At finite
length the gluing identity is
$\ip{e^{-tH_f}u_p}{e^{-tH_f}u_q}
=\ip{u_p}{e^{-2tH_f}u_q}$; the limiting pairing is the ordinary
positive harmonic Gram matrix.

\begin{proposition}\label{caps:descent}
The cap in \eqref{caps:heat} is cutoff independent and depends only
on the Jacobi class of $p\Omega$.
\end{proposition}
\begin{proof}
Differences of the cutoff operators in \eqref{caps:cutoff} are
$D_f$-commutators supported away from critical points, so their
action on a closed polynomial top form is compactly supported and
exact. If $p\Omega=df\wedge\gamma$ with $\gamma$ a polynomial
holomorphic one-form, then
$T_\chi(p\Omega)=D_f(T_\chi\gamma)$ is also compactly supported and
exact. Harmonic projection annihilates both kinds of difference.
The Hodge--Jacobi comparison identifies the resulting classes with
$\mathcal V_f$.
\end{proof}

The same assertion fails for the unmodified real-cycle integral.
For $f_c=G_c/4$ and $s=A^2+B^2$, $q_4=A^4+B^4+A^2B^2$, let
\[
z_c=Af_{c,A}+Bf_{c,B}=\frac{s+cq_4}{4}.
\]
This is zero in the Jacobi ring. Real integration by parts nevertheless
gives
\[
\frac1{2\pi}\int_{\R^2}z_c e^{-G_c}\,dA\,dB
   =\frac12 \mathcal C(c)>0\qquad(c>0).
\]
Indeed, integrating the divergence of $(A,B)e^{-G_c}$ gives
$\int (s+cq_4)e^{-G_c}=2\int e^{-G_c}$. Thus the raw period cannot
be identified with a physical cap functional on the same source
classes.

\subsection{Opposite twist and absolute normalization}
Let $U=(-1)^{\text{holomorphic degree}}$. It unitarily intertwines
$f$ and $-f$. Hodge star followed by complex conjugation gives an
antiunitary map $\mathfrak C:\mathcal V_{-f}\to\mathcal V_f$.
In middle degree on $\C^2$, $\mathfrak C^2=1$ and
$\mathfrak C U=U\mathfrak C$. Hence
$\Theta=\mathfrak C U$ is an antiunitary involution of $\mathcal V_f$.
The reflected caps and their functionals are
\begin{equation}\label{caps:reflection}
b_p=\Theta h_f(p),\qquad
\ell_p(q)=\ip{b_p}{h_f(q)}
 =\int_{\C^2}h_f(q)\wedge h_{-f}(p).
\end{equation}
The last pairing is complex bilinear. It is not a positive norm.
Opposite twists are essential: for a form $a$ of degree $k$,
\[
\int D_fa\wedge b=(-1)^{k+1}\int a\wedge D_{-f}b.
\]
This identity supplies cohomological descent of the wedge pairing.

The proportionality to the residue pairing is fixed by a Gaussian,
not chosen to match an Euler factor. In one coordinate, for
$f=az^2/2$ and $a\ne0$, representatives of $dz$ are
\[
\alpha_\pm=e^{-|a||z|^2}
 \left(dz\mp\frac{\bar a}{|a|}\,d\bar z\right).
\]
Their exact class normalization follows by applying $D_{\pm f}$ to
$(e^{-|a||z|^2}-1)/(\pm az)$. Direct integration gives
$\norm{\alpha_\pm}^2=2\pi/|a|$ and
$\int\alpha_+\wedge\alpha_-=-2\pi i/a$.
The additional reordering sign for two coordinates yields
$4\pi^2/(ab)$ for a diagonal Hessian $(a,b)$. With the residue
orientation fixed by $\Omega$, localization and deformation of the
cohomological pairing therefore give
\begin{equation}\label{caps:residue}
\ell_p(q)=4\pi^2\operatorname{Res}_f(pq).
\end{equation}
The opposite-twist Hodge pairing underlying this calculation is
discussed in \cite{LiWen}; the displayed constant is specific to the
Euclidean and differential conventions used here.

In the shape frame $E$ of \eqref{shape:frame}, put
\[
\lambda=\frac{256\pi^2}{3},\qquad
P=4\pi^2\eta_E
 =\lambda
\begin{pmatrix}
0&0&1\\0&1&-4/3\\1&-4/3&4/9
\end{pmatrix}.
\]
If $G_{ij}=\ip{h_i}{h_j}$ is the physical Gram matrix and
$\mathcal B_{ij}=\ip{b_i}{b_j}$ is the cap Gram matrix, then
\begin{equation}\label{caps:reality}
\mathcal B=P G^{-1}P^\dagger=G^T=\bar G.
\end{equation}
The first equality is the Riesz representation formula; the second
uses the antiunitary involution. In particular, residue data alone
do not determine $G$.

\subsection{The exact boundary Gram identity}
For any positive finite-dimensional metric $G$ and genuine boundary
map $v\mapsto y=Pv$, let $\mathcal B=PG^{-1}P^\dagger$. On the
range of $P$, the minimum-norm state with boundary values $y$ is
$v_{\min}=G^{-1}P^\dagger\mathcal B^+y$, where $+$ denotes the
Moore--Penrose inverse. Orthogonal decomposition gives
\begin{equation}\label{caps:gram}
\norm v_G^2=y^\dagger\mathcal B^+y+
             \norm{v-v_{\min}}_G^2.
\end{equation}
For one functional $\ell(v)=pv$ this reads
$\norm v_G^2=|\ell(v)|^2/(pG^{-1}p^\dagger)+\norm{v_\perp}_G^2$.
For invertible $P$ all states are detected and
$G=P^\dagger\mathcal B^{-1}P$. This is the physical replacement for
equating one raw period to one norm. For example,
$\ell_1(1)=0$ although $\norm{h_f(1)}^2>0$; the top cap detects
the identity through $\ell_v(1)=\lambda$.

\subsection{A normalized homogeneous prime source}
At $f_0=(Z^4+W^4+Z^2W^2)/16$, rotation symmetry makes the physical
metric in $(1,s_0,t_0)\Omega_Z$, where
$s_0=Z^2+W^2$ and $t_0=Z^2W^2$, diagonal. Its residue matrix is
$\lambda$ times the anti-diagonal identity. Equation
\eqref{caps:reality} then gives
\begin{equation}\label{caps:homogeneous}
G_0=\operatorname{diag}(a_0,\lambda,\lambda^2/a_0),
\qquad a_0>0.
\end{equation}
The middle class is self-reflected after normalization:
$\chi=h_{f_0}(s_0)/\sqrt\lambda$ has norm and bilinear self-pairing
equal to one.

For the homogeneous prime model
$f_M=M^2(U^4+V^4+U^2V^2)/16$, $M\ne0$, the change of fields
$\phi_M(U,V)=(\sqrt M\,U,\sqrt M\,V)$ gives the source
\begin{equation}\label{caps:prime-source}
\sigma_M=\frac{\sqrt3 M}{16\pi}(U^2+V^2)\,dU\wedge dV
 =\frac{1}{M\sqrt\lambda}\phi_M^*(s_0\Omega_Z).
\end{equation}
Degree-two Euclidean norms are invariant under the relevant
dilation. Thus its harmonic norm is $|M|^{-2}$ and its bilinear
self-pairing is $M^{-2}$. Taking $M=1-q$ gives a normalized
interacting Euler metric. This uses the middle Jacobi class,
not the identity class. The mass direction is still a dilation
and remains zero in the Jacobi ring.
